%% =====================================================================================
%%  Paper A -- Systematic review and critical appraisal (PRISMA 2020 structure)
%%  Methodological quality and reporting completeness of machine learning models for
%%  mental health prediction: a systematic review and PROBAST+AI / TRIPOD+AI appraisal
%%
%%  STATUS: Methods complete and protocol-grade. Results sections are structured
%%  placeholders (marked \todo) to be filled from the extraction workbook once
%%  screening and dual extraction are done. Numbers from the empirical companion
%%  (Paper B, this repository) are already in place in the Discussion.
%%  Compile: pdflatex -> bibtex not needed (thebibliography) -> pdflatex x2
%% =====================================================================================
\documentclass[11pt,a4paper]{article}
\usepackage[T1]{fontenc}\usepackage[utf8]{inputenc}\usepackage{lmodern}\usepackage{microtype}
\usepackage[margin=2.4cm]{geometry}\usepackage{booktabs}\usepackage{tabularx}\usepackage{array}
\usepackage{ragged2e}\usepackage{graphicx}\usepackage{xcolor}\usepackage{enumitem}
\usepackage[hidelinks]{hyperref}\usepackage{lineno}
% PRISMA 2020 flow diagram skeleton. Fill the macros from search_log.csv and Rayyan exports;
% the arithmetic is checked at compile time so that the figure cannot disagree with the text.
\newcommand{\nDb}{0}          % records identified from databases
\newcommand{\nOther}{0}       % records identified by citation chasing
\newcommand{\nDup}{0}         % duplicates removed
\newcommand{\nScreened}{\the\numexpr\nDb+\nOther-\nDup\relax}
\newcommand{\nExclTA}{0}      % excluded at title/abstract
\newcommand{\nSought}{\the\numexpr\nScreened-\nExclTA\relax}
\newcommand{\nNotRetrieved}{0}
\newcommand{\nAssessed}{\the\numexpr\nSought-\nNotRetrieved\relax}
\newcommand{\nExclFT}{0}      % excluded at full text (sum of reasons below)
\newcommand{\nIncluded}{\the\numexpr\nAssessed-\nExclFT\relax}
% Usage in the manuscript: \nIncluded{} studies were included (Figure~\ref{fig:prisma}).

\newcolumntype{L}{>{\RaggedRight\arraybackslash}X}
\newcommand{\todo}[1]{\textcolor{red}{[TODO: #1]}}
\setlength{\parskip}{0.4em}\setlength{\parindent}{0pt}
\title{Methodological quality and reporting completeness of machine learning models for mental health prediction: a systematic review and PROBAST+AI / TRIPOD+AI appraisal}
\author{Aminat Abiola Ajibola$^{1,*}$ \and Roger Nick Anaedevha$^{2}$}
\date{\footnotesize $^{1}$College of Computer Science and Engineering, University of Hafr Al-Batin, Eastern Province, Kingdom of Saudi Arabia \\
$^{2}$Institute of Cyber Intelligence Systems, National Research Nuclear University MEPhI, Moscow, Russian Federation \\[0.3em]
$^{*}$Corresponding author: aajibola@uhb.edu.sa \quad (R.N.A.: roger@robustidps.ai) \\[0.4em]
Protocol registered on PROSPERO: \todo{CRD number}}
\begin{document}
\maketitle
\linenumbers

\begin{abstract}
\textbf{Background.} Machine learning (ML) models that predict, detect or classify mental health
conditions are published at an accelerating rate, and descriptive reviews of the field already
exist. Whether these models are developed and reported to the standards now expected of clinical
prediction models has not been assessed across populations, disorders and data modalities.
\textbf{Methods.} We searched PubMed, Scopus, Web of Science and IEEE Xplore (1 January 2018 to
\todo{search date}) for studies developing or validating a supervised ML model predicting a mental
health outcome in humans, with at least one performance metric. Two reviewers screened,
extracted and appraised every study with PROBAST+AI (risk of bias and applicability) and the
TRIPOD+AI checklist (reporting completeness). Pre-specified secondary outcomes were the
proportions of studies with external validation, calibration assessment, discrimination reported
as AUROC rather than accuracy alone, leakage-prone evaluation design, events per variable (EPV)
relative to current guidance, and code or data availability.
\textbf{Results.} \todo{n records; n included; PROBAST+AI domain distribution; TRIPOD+AI
adherence; secondary-outcome proportions with 95\% CIs; stratification by modality.}
\textbf{Conclusions.} \todo{one sentence on prevalence of the deficiencies}. A companion empirical
study on three public actigraphy cohorts quantifies what these deficiencies cost: record-wise
cross-validation inflated AUROC by 0.04 to 0.17, external validation reduced AUROC by up to 0.21
with calibration slopes collapsing to 0.3, and no published model on those cohorts reports
calibration at all.
\end{abstract}

\section{Introduction}
Prediction models based on machine learning are proposed for nearly every mental health
condition and data source: electronic health records, questionnaires, wearable actigraphy, speech,
neuroimaging and social media text. Descriptive reviews catalogue the algorithms and datasets in
use \cite{hasan2026,madububambachu2024}; the most recent screened 3{,}320 records and organised
the field by method and data type \cite{hasan2026}. What none of them asks is whether the models
are trustworthy: whether they were evaluated on participants unseen in training, whether their
probabilities are calibrated, whether their sample sizes support the number of candidate
predictors, and whether the reports contain what a reader needs to reproduce or appraise them.

For clinical prediction models generally, that question has been asked and the answer is
discouraging. Meehan and colleagues appraised 308 psychiatric prediction models and found that
22.1\% tested calibration and 11.4\% performed internal validation to an adequate standard
\cite{meehan2022}; external validation is rare and generalisability poorly characterised
\cite{molpsych2025}. Those appraisals concentrate on regression-based clinical models. The ML
literature that now dominates mental-health prediction, and in particular the sensor-, speech-
and text-based studies, has been appraised only in narrow slices: adolescent EHR crisis prediction
\cite{adolescentehr2025}, psychosis prediction in at-risk individuals \cite{psychosisarms}, ADHD
\cite{adhd2024}, and social-media detection \cite{socialmedia2024}. Two developments make a
field-wide appraisal both feasible and overdue: TRIPOD+AI \cite{collins2024} and PROBAST+AI
\cite{moons2025} extend the reporting and risk-of-bias instruments explicitly to ML models,
adding items on data leakage, class imbalance, fairness, and ML-specific sample-size expectations.

Two evaluation practices are specific enough to ML on repeated-measures data to deserve their own
outcome. First, \emph{record-wise} splitting, in which windows or sessions from one participant
appear on both sides of a train/test split, lets a flexible model recognise the participant rather
than the condition; the resulting optimism is well documented in wearable and neuroimaging
research \cite{saeb2017,jmirai2026} but its prevalence in mental-health ML has not been counted.
Second, resampling or feature selection applied before splitting leaks test information into
training. Both are invisible in a results table and both are what turns a 0.75 AUROC into a 0.95.

We therefore asked: how methodologically sound and completely reported are ML models for mental
health prediction, across all populations, disorders and modalities? We pre-specified the
proportions of studies affected by each of six deficiencies and, in a paired empirical study on
public actigraphy cohorts \cite{paperB}, measured how much each of the three most consequential
ones costs.

\section{Methods}
The review follows PRISMA 2020 \cite{page2021} and was registered on PROSPERO (\todo{CRD}). The
full protocol, search strings, extraction workbook and code are available in the project
repository \cite{repo}.

\subsection{Eligibility criteria}
\begin{table}[!htbp]\centering\caption{Inclusion and exclusion criteria.}\label{tab:criteria}\footnotesize
\begin{tabularx}{\linewidth}{@{}p{2.2cm} L L@{}}\toprule
\textbf{Element} & \textbf{Inclusion} & \textbf{Exclusion}\\\midrule
Population & Human participants, any age or setting, any country & Animal studies; simulation-only work\\
Condition & Any mental health disorder or state (depressive, anxiety, stress-related, PTSD, psychotic, bipolar, eating, substance use, ADHD, psychological distress) & Neurological conditions without a psychiatric outcome; physical health outcomes\\
Index model & Supervised ML or DL model developed and/or validated to predict, detect, classify or diagnose a mental health outcome & Descriptive statistics; unsupervised clustering without a predictive target; conventional regression alone; chatbots or interventions without a prediction model\\
Comparator & Not required & ---\\
Outcomes & At least one quantitative performance metric & No performance metric\\
Design & Development, validation, or development-with-validation studies; any data source & Reviews, editorials, protocols, abstracts without full text, preprints (analysed as a separate stratum if included)\\
Timeframe & 1 January 2018 to \todo{date} & Before 2018\\
Language & English & Non-English\\
\bottomrule\end{tabularx}\end{table}

\subsection{Information sources and search}
PubMed, Scopus, Web of Science Core Collection and IEEE Xplore were searched on \todo{date}
using three concept blocks (machine learning AND mental health AND prediction/detection); the
full strings are in Supplementary Table S1 and the repository. PsycINFO \todo{was / was not}
searched \todo{state reason}. Backward and forward citation chasing was performed for every
included study and for the four overlapping reviews \cite{hasan2026,meehan2022,molpsych2025,madububambachu2024}.
\todo{Record counts per database from search\_log.csv.}

\subsection{Selection and data collection}
Records were de-duplicated and screened in Rayyan with blinding by two reviewers independently at
title/abstract and full-text stages; disagreements were resolved by discussion or a third
reviewer. Inter-rater agreement is reported as Cohen's $\kappa$. Data were extracted in duplicate
into a structured workbook (Supplementary File S2) after a five-study pilot: study characteristics,
population, outcome definition and prevalence, data modality, number of candidate predictors,
models, validation design, performance metrics, and availability statements.

\subsection{Risk of bias and reporting appraisal}
Risk of bias and applicability were assessed with PROBAST+AI \cite{moons2025} by two reviewers
independently (domains: participants and data, predictors, outcome, analysis). Reporting
completeness was scored against the 27 TRIPOD+AI items \cite{collins2024} as fully, partially or
not reported, with adherence per study computed as fully reported items over applicable items.

\subsection{Outcomes}
Primary: PROBAST+AI risk-of-bias judgement per domain and overall; TRIPOD+AI adherence per item and
per study. Secondary (pre-specified): (1) any external validation (temporal, geographical or
independent cohort; a random hold-out does not qualify); (2) calibration reported by any method
(plot, slope/intercept, Brier, ECE, Hosmer--Lemeshow); (3) AUROC reported versus accuracy as the
only discrimination metric; (4) leakage-prone evaluation design, coded as present when
repeated-measures data were split record-wise, when resampling or feature selection preceded the
split, or when a predictor was derived from the outcome; (5) EPV, computed as outcome events
divided by candidate predictors and compared with the $\geq$20 (regression) and higher ML
expectations discussed in PROBAST+AI \cite{moons2025}; (6) code and data availability with a
working link. Studies reporting $\geq$99\% accuracy without external validation were examined
individually as leakage case studies.

\subsection{Synthesis}
Proportions with Wilson 95\% confidence intervals, overall and stratified by data modality
(EHR/registry, questionnaire, wearable/actigraphy, speech/interview, imaging, social-media text),
disorder and publication year. Meta-analysis of performance was pre-specified only for
clinically homogeneous subsets reporting AUROC with variance; \todo{state whether performed}.

\section{Results}
\subsection{Study selection}
\todo{PRISMA flow: \nDb{} database records, \nOther{} from other sources, \nDup{} duplicates,
\nScreened{} screened, \nExclTA{} excluded, \nAssessed{} full texts assessed, \nIncluded{} included.
Use paper/review/prisma\_flow\_template.tex.}
\subsection{Study characteristics}
\todo{Table: year, country, population, disorder, modality, sample size, prevalence, models. Summarise medians and IQRs.}
\subsection{Risk of bias (PROBAST+AI)}
\todo{Stacked bar per domain; proportion high/unclear/low; which signalling questions drive high risk (expected: analysis domain, 4.1 sample size, 4.7 calibration, 4.8 optimism/leakage).}
\subsection{Reporting completeness (TRIPOD+AI)}
\todo{Per-item adherence bar chart; median adherence per study; items with lowest adherence (expected: sample size justification, calibration, fairness, code availability, model specification).}
\subsection{Secondary outcomes}
\todo{Table with one row per outcome: n/N (\%, 95\% CI), stratified by modality. Leakage case studies described narratively.}

\section{Discussion}
\subsection{Principal findings}
\todo{Three sentences on the headline proportions.}

\subsection{What the deficiencies cost: the companion study}
The proportions above describe how widespread each practice is; they do not say how much it
matters. The companion study \cite{paperB} answers that on three public actigraphy cohorts
(DEPRESJON, PSYKOSE, HYPERAKTIV; 162 unique participants, 1{,}818 unique valid day-windows, the 32 healthy controls being shared by the first two cohorts) with the same
fixed models, seeds and subject-level bootstrap. Record-wise cross-validation inflated window-level
AUROC by 0.065 (95\% CI 0.024 to 0.122) to 0.114 (0.060 to 0.181) for depression versus controls,
by 0.020 to 0.044 for schizophrenia versus controls, and by 0.130 (0.095 to 0.167) to 0.169
(0.135 to 0.206) for ADHD versus clinical controls, where the honest estimate was at chance
(AUROC 0.44 to 0.45): leakage did not merely exaggerate a real signal, it manufactured one.
Applying a model trained on one cohort to another without refitting reduced AUROC by 0.13 to 0.21
in the harder direction and collapsed calibration slopes to 0.27 to 0.63 with intercepts of
0.7 to 1.4, i.e.\ probabilities were both overconfident and shifted, while a deliberate label-shift
transfer (depression model applied to ADHD versus clinical controls) fell to chance. Even under
honest internal validation, calibration slopes were 0.53 to 0.95 and expected calibration error
0.06 to 0.22, and no published model on these cohorts reports any of these quantities. Subject-level
events per candidate predictor were 0.6 to 0.8 in every cohort, one to two orders of magnitude
below guidance, although window-level counts make the samples look ten to twenty times larger.

\subsection{Comparison with previous reviews}
Our findings extend Meehan et al.\ \cite{meehan2022} from regression-based clinical models to the
ML and digital-phenotyping literature, and go beyond the descriptive catalogue of Hasan et al.\
\cite{hasan2026} by appraising quality rather than method. Benchmark-level audits of clinical
interview corpora \cite{daicaudit2026,daicinvalid2025} reach the same conclusion by a different
route: published rankings on DAIC-WOZ are unstable across splits and models learn corpus artefacts.

\subsection{Strengths and limitations}
Breadth across disorders and modalities; pre-registration; PROBAST+AI and TRIPOD+AI applied in
duplicate; a paired empirical companion. Limitations: English-only; \todo{PsycINFO}; appraisal
depends on what authors report, so under-reporting and poor conduct cannot always be separated;
\todo{single second screener on a sample, if applicable}.

\subsection{Implications}
Journals in this field should require, as a condition of review, a statement of the unit of
splitting for repeated-measures data, a calibration plot or slope, subject-level (not window-level)
sample sizes with EPV, and a code link. Authors should treat the TRIPOD+AI checklist as a design
document rather than a submission form.

\section*{Declarations}
Funding: \todo{}. Conflicts of interest: none declared. Registration: PROSPERO \todo{CRD}. Data
and code: extraction workbook, search log and analysis code at \cite{repo}.

\begin{thebibliography}{99}\footnotesize
\bibitem{hasan2026} Hasan MJ, Shifat SH, Matubber J, et al. An in-depth exploration of machine learning methods for mental health state detection: a systematic review and analysis. Front Digit Health. 2026. doi:10.3389/fdgth.2025.1724348
\bibitem{madububambachu2024} Madububambachu U, et al. Machine learning techniques to predict mental health diagnoses: a systematic literature review. Clin Pract Epidemiol Ment Health. 2024.
\bibitem{meehan2022} Meehan AJ, Lewis SJ, Fazel S, et al. Clinical prediction models in psychiatry: a systematic review of two decades of progress and challenges. Mol Psychiatry. 2022;27:2700--2708. doi:10.1038/s41380-022-01528-4
\bibitem{molpsych2025} Generalizability of clinical prediction models in mental health. Mol Psychiatry. 2025. doi:10.1038/s41380-025-02950-0
\bibitem{adolescentehr2025} Machine learning models for predicting mental health crises in adolescents using electronic health records: a systematic review. 2025. PMC12426581.
\bibitem{psychosisarms} Systematic review of clinical prediction models for psychosis in individuals meeting At Risk Mental State criteria. PMC11480010.
\bibitem{adhd2024} Individualized prediction models in ADHD: a systematic review and meta-regression. Mol Psychiatry. 2024. doi:10.1038/s41380-024-02606-5
\bibitem{socialmedia2024} Machine learning approaches for mental illness detection on social media: a systematic review of biases and methodological challenges. arXiv:2410.16204.
\bibitem{collins2024} Collins GS, Moons KGM, Dhiman P, et al. TRIPOD+AI statement: updated guidance for reporting clinical prediction models that use regression or machine learning methods. BMJ. 2024;385:e078378.
\bibitem{moons2025} Moons KGM, Damen JAA, Nair T, et al. PROBAST+AI: an updated quality, risk of bias, and applicability assessment tool for prediction models using regression or artificial intelligence methods. BMJ. 2025;388:e082505.
\bibitem{saeb2017} Saeb S, Lonini L, Jayaraman A, Mohr DC, Kording KP. The need to approximate the use-case in clinical machine learning. GigaScience. 2017;6:1--9.
\bibitem{jmirai2026} Karbalaie A, Abtahi F, H\"ager C. Participant-aware model validation for repeated-measures data: comparative cross-validation study. JMIR AI. 2026;5:e87728.
\bibitem{page2021} Page MJ, McKenzie JE, Bossuyt PM, et al. The PRISMA 2020 statement. BMJ. 2021;372:n71.
\bibitem{daicaudit2026} A multi-probe audit of clinical-interview depression detection benchmarks. arXiv:2605.23977. 2026.
\bibitem{daicinvalid2025} Most DAIC-WOZ depression classifiers are invalid, they don't learn task-specific features: preliminary findings from a large-scale reproducibility study. ICMI Companion. 2025. doi:10.1145/3747327.3763034
\bibitem{paperB} Ajibola AA, et al. How much do leaky splits, missing external validation and missing calibration cost? An empirical companion on public actigraphy cohorts. \todo{preprint DOI}. 2026.
\bibitem{repo} mlmh-appraisal repository. \url{https://github.com/rogerpanel/cv/tree/main/mlmh-appraisal}. \todo{tag and DOI via Zenodo}.
\end{thebibliography}
\end{document}
